\DocumentMetadata{
  pdfversion=2.0,
  pdfstandard=ua-2,
  lang=en-US,
  tagging=on,
  tagging-setup={math/setup=mathml-SE,tabsorder=structure}
}
\documentclass[11pt,letterpaper]{article}
\usepackage[margin=0.68in,includefoot]{geometry}
\usepackage{newcomputermodern}
\usepackage{amsmath,amscd,mathtools}
\usepackage{tikz,adjustbox}
\providecommand{\square}{\mdlgwhtsquare}
\usepackage[hidelinks]{hyperref}
\hypersetup{
  pdftitle={Numerical Analysis PhD Exam, May 2012},
  pdfauthor={University of Florida Department of Mathematics},
  pdfsubject={Graduate mathematics examination},
  pdfdisplaydoctitle=true,
  bookmarksopen=true
}
\setcounter{secnumdepth}{0}
\setlength{\parindent}{0pt}
\setlength{\parskip}{0.28em}
\setlength{\emergencystretch}{3em}
\setlength{\leftmargini}{2.15em}
\setlength{\leftmarginii}{2.65em}
\setlength{\leftmarginiii}{3em}
\renewcommand{\labelenumi}{\textbf{\theenumi.}}
\pagestyle{empty}
\raggedbottom
\allowdisplaybreaks
\newenvironment{examproblems}[1]{%
  \begin{enumerate}%
  \setcounter{enumi}{#1}%
  \setlength{\topsep}{0.35em}%
  \setlength{\partopsep}{0pt}%
  \setlength{\itemsep}{0.48em}%
  \setlength{\parsep}{0.18em}%
}{\end{enumerate}}
\newenvironment{examsubparts}{%
  \begin{enumerate}%
  \setlength{\topsep}{0.18em}%
  \setlength{\partopsep}{0pt}%
  \setlength{\itemsep}{0.16em}%
  \setlength{\parsep}{0.1em}%
}{\end{enumerate}}
\newenvironment{examinstructions}{%
  \par\begingroup\itshape
}{%
  \par\endgroup\vspace{0.18em}
}
\makeatletter
\renewcommand\section{\@startsection{section}{1}{0pt}%
  {-1.25ex plus -.25ex minus -.1ex}%
  {0.55ex plus .1ex}%
  {\normalfont\large\bfseries}}
\renewcommand{\@maketitle}{%
  \newpage\null\vskip -1em%
  {\footnotesize\bfseries
    \href{https://www.ufl.edu/}{University of Florida}, \href{https://math.ufl.edu/}{Department of Mathematics}\par}%
  \vskip -0.5em%
  \tagstructbegin{tag=Title}%
    {\LARGE\bfseries\href{https://gma.math.ufl.edu/exams/numerical-analysis-phd/}{Numerical Analysis PhD Exam}, May 2012\par}%
  \tagstructend%
  \vskip 0.25em%
  \tagmcbegin{artifact}%
    \hrule height 1pt%
  \tagmcend%
  \vskip 1em%
}
\makeatother
\title{Numerical Analysis PhD Exam, May 2012}
\author{}
\date{}
\begin{document}
\maketitle
\thispagestyle{empty}
\begin{examinstructions}
Answer any 8 questions.
\end{examinstructions}
\begin{examproblems}{0}
\item
Suppose that~\(A\) and \(B\in\mathbb{C}^{n\times n}\) are Hermitian matrices. If \(\sigma(A)\) denotes the largest eigenvalue of~\(A\), then show that 
\[
\sigma(A+B)\leq \sigma(A)+\sigma(B).
\]
\item
Suppose that \(A\in\mathbb{C}^{n\times n}\) is an invertible matrix, \(u\in\mathbb{C}^n\), and \(v\in\mathbb{C}^n\). Let \(v^*\) denote the conjugate transpose of~\(v\).
\begin{examsubparts}
\item[\textnormal{(a)}]
Show that 
\[
\det(A-uv^*)=(1-v^*A^{-1}u)\det A.
\]
\item[\textnormal{(b)}]
Show that \(A-uv^*\) is invertible if and only if \(v^*A^{-1}u\neq 1\). Moreover, if \(v^*A^{-1}u\neq 1\), then 
\[
(A-uv^*)^{-1}=A^{-1}+\left(\frac{1}{1-v^*A^{-1}u}\right)A^{-1}uv^*A^{-1}.
\]
\end{examsubparts}
\item
This problem has two parts.
\begin{examsubparts}
\item[\textnormal{(a)}]
Suppose~\(p\) and \(q\in\mathbb{R}\) with~\(p\) and~\(q\) positive and \(p^{-1}+q^{-1}=1\). Show that for any matrix \(A\in\mathbb{C}^{m\times n}\), we have \(\lVert A\rVert_p=\lVert A^*\rVert_q\).
\item[\textnormal{(b)}]
Prove that 
\[
\lVert A\rVert_2^2\leq \lVert A\rVert_p\lVert A\rVert_q
\]
 for any \(A\in\mathbb{C}^{m\times n}\) and any positive~\(p\) and \(q\in\mathbb{R}\) with \(p^{-1}+q^{-1}=1\).
\end{examsubparts}
\item
This problem has two parts.
\begin{examsubparts}
\item[\textnormal{(a)}]
For any matrices~\(A\) and~\(B\), show that the nonzero eigenvalues of \(AB\) and \(BA\) are the same.
\item[\textnormal{(b)}]
If \(AB\) is normal, \(\lVert\cdot\rVert_2\) is the 2-norm of a matrix, and \(\lVert\cdot\rVert\) is an induced matrix norm, then show that \(\lVert AB\rVert_2\leq\lVert BA\rVert\).
\end{examsubparts}
\item
Let~\(P\) and~\(Q\) be two \(m\times m\) orthogonal projectors. Prove that \(\lVert P-Q\rVert_2\leq 1\).
\item
Consider the following minimization problem: 
\[
\tau_n=\inf_{\deg Q<n}\left(\max_{x\in[-1,1]}\lvert x^n+Q(x)\rvert\right).
\]
 Prove that \(\tau_n=\frac{1}{2^{n-1}}\), the infimum is in fact a minimum, attained at a unique \(Q^*\) satisfying 
\[
x^n+Q^*(x)=\frac{1}{2^{n-1}}T_n(x),
\]
 where \(\{T_n\}\) are the Chebyshev’s polynomials.
\item
Design an algorithm with a cubic rate of convergence for computing the quantity \(\sqrt{5}\). Prove that your algorithm is indeed cubic. Find an interval \([a,b]\subset[0,+\infty)\) such that any iterative sequence starting in \([a,b]\) will converge to \(\sqrt{5}\).
\item
State the classical Hermite interpolation problem. Prove that the problem is well-posed, i.e. prove the existence and uniqueness of solutions. Derive the error formula for the interpolating polynomial.
\item
Describe Simpson’s Rule for numerical integration. State (without proof) the error formula of the Trapezoidal Rule. Show that the error bound cannot be improved.
\item
Let \(w(x)>0\) be integrable on \([a,b]\). Define an orthogonal polynomial family \(\{\phi_n\}\) on \([a,b]\) with weight \(w(x)\). Prove that between two consecutive zeros of \(\phi_n\) there is exactly one zero of \(\phi_{n-1}\). You may use the triple recursion formula without proof.
\end{examproblems}
\end{document}
